% ID: SEL-D0025-OB-E04
% ESTADO: cotejada. Segunda pasada visual completa y borrador compilado revisados.
% Fuente: Selectividad/matematicas_aplicadas/suplemente_junio_EXAMEN 4-ANDALUCÍA.pdf
% SHA-256: c1a3a7a37d38b04359115c64762b1f8a543e4d66155338931de84ef6c304da2c
% Página: 2; opción B; ejercicio original 4.

El peso de los paquetes de levadura de una marca sigue una ley Normal de desviación típica 0.3 g. Se desea construir un intervalo de confianza, al 98\%, para estimar la media. Para ello, se toma una muestra aleatoria de 9 paquetes.
\begin{enumerate}[label=\alph*)]
\item \textbf{(1.25 puntos)} \textquestiondown{}Qué amplitud tendrá dicho intervalo?
\item \textbf{(1.25 puntos)} Obtenga el intervalo sabiendo que los pesos, en gramos, de los paquetes son:
\[
10\quad9.9\quad10.04\quad9.5\quad10.1\quad9.8\quad10.2\quad10\quad10.3
\]
\end{enumerate}
